\documentclass[../BTL.tex]{subfiles}
\pagenumbering{roman}
\begin{document}
\begin{center}
    \Large{\textbf{LỜI CẢM ƠN}}\\
\end{center}
\vspace{1cm}
Nhóm em xin gửi lời cảm ơn tới thầy TS Phạm Hoàng Việt đã hướng dẫn nhóm em đề tài này, có những nhận xét và gợi ý trong mỗi buổi họp hàng tuần để giúp nhóm chúng em có thêm những ý tưởng và chỉnh sửa để hoàn thiện đề tài.
Nhóm em cũng xin gửi lời cảm ơn chân thành tới các thầy, cô trong Khoa Công nghệ thông tin 1 đã dạy những kiến thức cơ sở của ngành và chuyên ngành cho chúng em trong những năm tháng tại trường.
Ngoài ra chúng em cũng xin gửi lời tri ân tới gia đình, bạn bè đã luôn ở bên cạnh để động viên và hỗ trợ trong quá trình nhóm làm đề tài này. 


Với cá nhân em, Vũ Công Tuấn Dương, việc được học những người thầy tâm huyết, đặc biệt trong những năm tháng tại trường góp phần không nhỏ trong việc hình thành cách làm việc, cách suy nghĩ của em. Em xin gửi những lời cảm ơn chân thành đến những người thầy để lại trong em nhiều sự ngưỡng mộ và ấn tượng nhất trong quá trình học : Em may mắn được học thầy Nguyễn Xuân Anh môn Mạng máy tính, thầy là một người rất kỷ luật và nghiêm khắc nhưng rất quan tâm tới sinh viên. Học thầy, em học được tính tự giác, kỷ luật đôi khi đến mức cực đoan của thầy cũng như cách thầy tư duy về sự cần thiết có những tính toán dự phòng cho những trường hợp khẩn cấp, tư duy về việc tối ưu công việc, những quan điểm về hệ thống phần mềm em rút ra từ trang web và phần mềm thi của thầy. Bản thân thầy tự rèn luyện để chính cách thầy sống là tấm gương cho những tiêu chuẩn và nguyên tắc của thầy. Em vẫn luôn dành sự tôn trọng nhất đến thầy, em cảm ơn thầy rất nhiều. Người thầy tiếp theo mà em muốn dành lời cảm ơn, thầy Nguyễn Mạnh Hùng dạy em môn Nhập môn Công nghệ phần mềm. Trong suốt một học kỳ có cơ hội được học thầy, được làm việc với những yêu cầu khắt khe của thầy, em đã học được chút cách làm việc khoa học để có thể hoàn thành kịp hạn, hoàn thành tốt những bài tập thầy giao. Ngoài ra, em còn tự ngẫm và rút ra cho mình được một bài học từ những nội quy trong lớp của thầy, đó là nguyên tắc không làm phiền người khác, em áp dụng nguyên tắc này có thể giải thích dược các quy tắc ứng xử nơi công cộng. Thật sự trải nghiệm học môn của thầy khá vất vả nhưng vui và cuối kỳ nhận được điểm xứng đáng với những nỗ lực bỏ ra cho môn học của thầy, em rất cảm ơn thầy. Em xin dành lời tri ân tiếp theo tới thầy Nguyễn Mạnh Sơn, em có học thầy môn Lập trình hướng đối tượng, đây chính là môn học quan trọng, cho em nền tảng và cũng góp phần không nhỏ trong định hướng hiện tại và tương lai của em. Tại thời điểm học môn này - kỳ 1 năm thứ ba, em vẫn còn chưa có định hướng chuyên môn của mình sau này, khá hoang mang, nhưng sau khi học, làm bài tập lớn môn của thầy, em thấy hứng thú hơn với ngôn ngữ lập trình Java, tư duy hướng đối tượng nên dần tìm hiểu thêm và dần dần tìm, hình dung được con đường mà mình sẽ đi. Thầy thực sự là một người thầy tâm huyết với nghề, với sinh viên. Em vẫn nhớ những quan điểm của thầy về tương lai ngành, về những gì sinh viên cần chuẩn bị với tương lai phát triển rất nhanh của AI - đó là tìm hiểu thật kỹ về nghiệp vụ sản phẩm, theo thực tế. Đây là những phần mà AI chưa có làm tốt tại thời điểm thầy nói(2024), và thầy có tầm nhìn và quan điểm như vậy. Tính đến thời điểm hiện tại, em thấy quan điểm của thầy vẫn đúng - với sự hỗ trợ của AI, rất nhiều các sản phẩm phần mềm được làm ra với tốc độ rất nhanh, nhưng thực tế ứng dụng thì chưa đáp ứng đủ nhiều các nhu cầu người dùng. Em xin dành lời tri ân tiếp theo tới thầy Ngô Xuân Bách, em có cơ hội học thầy môn Toán rời rạc 1. Mặc dù hiện nay thầy không còn công tác, giảng dạy tại trường, nhưng em vẫn luôn ấn tượng với những bài giảng hay và giải thích chi tiết, tận gốc vấn đề, và rất dễ hiểu của thầy. Thầy là một trong số ít những người thầy mà theo em, giỏi và có thành tựu ở cả 2 mảng là doanh nghiệp và học thuật. Người thầy cuối cùng trong danh sách 5 người thầy em ấn tượng nhất trong thời gian học là thầy Đinh Chí Hiếu dạy môn Lập trình với Python. Tính đến thời điểm hiện tại, thầy là người duy nhất cho phép và khuyến khích sinh viên dùng tiếng Anh trong lớp. Đó là ấn tượng mà em không quên khi nhớ lại những ngày tháng học môn của thầy. Ngoài ra, những quan điểm, góc nhìn của một người từng có nhiều cơ hội làm việc với những doanh nghiệp, những người nước ngoài đã giúp em học hỏi và biết thêm nhiều điều từ thầy. Một trong những bài học, trải nghiệm mà em học được từ thầy là trong quá trình thuyết trình bài tập nhóm, mỗi người chỉ có từ 2 đến 3 phút để làm ấn tượng người nghe, phải trình bày tóm gọn và đầy đủ, chỉ trong thời gian như vậy là có thể biết được người nào giỏi và ngược lại. Thực sự đó là một trong những trải nghiệm vui và nhớ nhất của em đến tận bây giờ. Cảm ơn thầy rất nhiều


\end{document}